\documentclass[11pt,reqno]{amsart}

\usepackage[T1]{fontenc}
\usepackage{lmodern}
\usepackage{microtype}
\usepackage{mathtools}
\usepackage{amssymb}
\usepackage{mathrsfs}
\usepackage{booktabs}
\usepackage{enumitem}
\usepackage{xcolor}
\usepackage[margin=1.08in]{geometry}
\usepackage[colorlinks=true,linkcolor=blue!45!black,citecolor=blue!45!black,urlcolor=blue!45!black]{hyperref}
\hypersetup{
  pdftitle={Twisted Eckardt Periods, Uniform Prill-Exceptional Covers, and a Prym Fixed-Part Obstruction},
  pdfauthor={FDmd233},
  pdfsubject={Cubic threefolds, Raynaud bundles, Prym monodromy, and Putman--Wieland},
  pdfkeywords={Eckardt point, F4, Raynaud bundle, Prill problem, Prym map, Putman--Wieland}
}

\numberwithin{equation}{section}
\setcounter{MaxMatrixCols}{12}
\setlist[enumerate]{leftmargin=2.2em,itemsep=2pt,topsep=4pt}
\setlist[itemize]{leftmargin=2.2em,itemsep=2pt,topsep=4pt}

\newtheorem{theorem}{Theorem}[section]
\newtheorem{proposition}[theorem]{Proposition}
\newtheorem{lemma}[theorem]{Lemma}
\newtheorem{corollary}[theorem]{Corollary}
\newtheorem{claim}[theorem]{Claim}
\theoremstyle{definition}
\newtheorem{definition}[theorem]{Definition}
\theoremstyle{remark}
\newtheorem{remark}[theorem]{Remark}

\newcommand{\C}{\mathbf C}
\newcommand{\Q}{\mathbf Q}
\newcommand{\Z}{\mathbf Z}
\newcommand{\OO}{\mathcal O}
\newcommand{\PP}{\mathbf P}
\newcommand{\cM}{\mathcal M}
\newcommand{\cA}{\mathcal A}
\newcommand{\cR}{\mathcal R}
\newcommand{\cV}{\mathcal V}
\newcommand{\cE}{\mathcal E}
\newcommand{\rk}{\operatorname{rk}}
\newcommand{\im}{\operatorname{im}}
\newcommand{\Alb}{\operatorname{Alb}}
\newcommand{\Pic}{\operatorname{Pic}}
\newcommand{\Sym}{\operatorname{Sym}}
\newcommand{\Hom}{\operatorname{Hom}}
\newcommand{\Spec}{\operatorname{Spec}}
\newcommand{\Stab}{\operatorname{Stab}}
\newcommand{\Gal}{\operatorname{Gal}}
\newcommand{\Mod}{\operatorname{Mod}}
\newcommand{\Sp}{\operatorname{Sp}}
\newcommand{\SU}{\operatorname{SU}}
\newcommand{\GL}{\operatorname{GL}}
\newcommand{\SL}{\operatorname{SL}}
\newcommand{\IC}{\operatorname{IC}}
\newcommand{\id}{\mathrm{id}}
\newcommand{\GGG}{\mathrm{GGG}}

\title[Twisted Eckardt periods and Prym obstructions]
{Twisted Eckardt Periods, Uniform Prill-Exceptional Covers,\\
and a Prym Fixed-Part Obstruction}
\author{FDmd233}
\date{13 July 2026}

\subjclass[2020]{Primary 14H30, 14J30; Secondary 14D07, 14H40, 14K12, 20G41, 32G20, 57K20}
\keywords{cubic threefold, Fano surface, Eckardt point, exceptional group $F_4$, twisted period map, Raynaud bundle, Prill's problem, Prym map, Putman--Wieland conjecture}

\begin{document}

\begin{abstract}
I relate two constructions with finite-unitary coefficients: the Eckardt quotient of the Fano surface of a cubic threefold, whose convolution Tannaka group is $F_4$, and Raynaud's bundles without theta divisors on curves.  For an explicit marked Eckardt cubic, I prove that the twisted period differential is injective on the full seven-dimensional marked Eckardt tangent space.  Its second Higgs iterate has rank between $12$ and $21$ on a nonempty Zariski-open set of characters; the upper bound is forced by the Eckardt reflection.  For torsion characters the horizontal Tate line splits off, and the complementary rank-$26$ summand has period rank $7$ at the distinguished complex embedding.

For every $g\ge2$, I then construct a single irreducible finite representation
$\rho_g:\pi_1(\Sigma_g)\to \SU(g^g)$ which is Prill-exceptional for every marked complex structure and every point.  If $G_g$ is its image, then
\[
1+g^{2g}\le |G_g|\le g^{3g}.
\]
This gives a fixed topological Galois cover with at least $1+g^g$ sections of every fibre divisor and, for $g\ge3$, disproves the generic-global-generation conjecture for unitary local systems.  Independently, every Prill-exceptional connected \emph{\'etale} cover has degree at least $g+2$, and every such Galois cover has degree at least $1+(g+1)^2$.

Finally, I prove that maximal Siegel rank at one point of a weight-one period map excludes every virtual constant subvariation.  Applied to the degree-$27$ and degree-$3$ Prym maps, this prevents flat classes on the smooth-Hurwitz parts of special Fano and Eckardt fibres from extending to finite-index constant factors.  These double-cover no-finite-orbit cases were already known; no new case of the general Putman--Wieland conjecture is claimed.  Neither the weight-two $F_4$ summand nor failure of generic global generation alone produces a finite monodromy orbit.
\end{abstract}

\maketitle

\section{Introduction}

The two starting points of this paper have the same numerical shadow but live in different Hodge-theoretic categories.  Let $F$ be the Fano surface of lines on a smooth cubic threefold with an Eckardt point, let $E\subset F$ be the elliptic curve parametrizing the lines through that point, and let
\[
A=\Alb(F),\qquad B=A/E,\qquad u:F\longrightarrow B.
\]
The quotient contracts $E$ and satisfies
\begin{equation}\label{eq:semismall-intro}
Ru_*\C_F[2]\simeq \delta_X\oplus\delta_0,
\qquad X=u(F),
\end{equation}
while the convolution Tannaka group of $\delta_X$ is the simple group $F_4$.  Under the resulting restriction of the Albert representation,
\begin{equation}\label{eq:branching-intro}
\mathbf{27}|_{F_4}=\mathbf 1\oplus\mathbf{26}.
\end{equation}
The first goal is to determine whether the $26$-dimensional term in \eqref{eq:branching-intro} varies maximally along the Eckardt locus, rather than merely identifying its fibrewise Tannaka symmetry.

The second starting point is Raynaud's rank-$g^g$ bundle on a genus-$g$ curve.  After a determinant-normalized twist, it is a stable degree-zero bundle $V$ with finite monodromy such that
\begin{equation}\label{eq:raynaud-intro}
H^0(C,V)=0,
\qquad H^0(C,V(p))\ne0\quad\text{for every }p\in C.
\end{equation}
Condition \eqref{eq:raynaud-intro} is precisely the nontrivial summand required for a fibre divisor on a finite \'etale cover to move.  A priori the finite representation depends on $C$.  The second goal is to remove that dependence.
I call a finite \'etale cover $f:X\to C$ \emph{Prill-exceptional} if
$h^0(X,\OO_X(f^{-1}(p)))\ge2$ for every $p\in C$; a finite representation is Prill-exceptional when its associated summand has the equivalent nonvanishing property in \eqref{eq:raynaud-intro}.

Here are the main results.  The precise forms, including coefficient fields and marked-moduli qualifications, are stated in the body of the paper.

\begin{theorem}[Eckardt twisted infinitesimal Torelli]\label{thm:intro-eckardt}
For the marked Eckardt pair
\[
Y_0=\{x_0^2x_1+x_1^3+x_2^3+x_3^3+x_4^3=0\}\subset\PP^4,
\qquad p_0=[1:0:0:0:0],
\]
let $T_E$ be the tangent space to the local moduli space of marked Eckardt pairs at $(Y_0,p_0)$.  Then $\dim T_E=7$.  There is a nonempty Zariski-open set
\[
\Omega\subset\Pic^0(B)\setminus\{0\}
\]
such that, for $M\in\Omega$ and $L=u^*M$, the infinitesimal Higgs map on $H^2(F,L)$ is injective on $T_E$.  Its second iterate satisfies
\[
12\le\rk\theta_L^{(2)}|_{\Sym^2T_E}\le21.
\]
The set $\Omega$ contains torsion points of arbitrarily large order.  For a torsion point, the horizontal Tate line in \eqref{eq:semismall-intro} splits off.  The complementary summand has Hodge numbers $(6,14,6)$ at every complex embedding of its cyclotomic field, and period rank $7$ along $T_E$ at the distinguished embedding defined by $M$.
\end{theorem}

The proof is an exact Jacobian-ring computation.  The image
\[
I_E=\mu(\Sym^2T_E)
\]
has dimension $27$.  The reflection identifies the two six-dimensional extreme Hodge spaces and makes the projected multiplication symmetric, which gives the global upper bound $\binom{6+1}{2}=21$.  A saturated rank-six limit at the trivial character produces a $12\times12$ determinant $-2^7 3^6$ and a flattened $7\times7$ determinant $-2^3 3^7$.  The latter proves injectivity of the first Higgs field by integrability.  The rank-six central space is the fibre of a direct-image bundle; it is not the ten-dimensional untwisted space $H^0(F,K_F)$.

\begin{theorem}[Uniform Raynaud--Prill representation]\label{thm:intro-uniform}
Let $g\ge2$ and $r=g^g$.  There is a nontrivial irreducible representation with finite image
\[
\rho_g:\pi_1(\Sigma_g,x_0)\longrightarrow \SU(r)
\]
such that, for every marked compact Riemann surface $C$ of genus $g$ and every $p\in C$, the associated bundle $V_{\rho_g,C}$ satisfies
\[
H^0(C,V_{\rho_g,C})=0,
\qquad H^0(C,V_{\rho_g,C}(p))\ne0.
\]
If $G_g=\im\rho_g$, then
\[
1+g^{2g}\le |G_g|\le g^{3g}.
\]
The fixed topological Galois cover associated with $\ker\rho_g$ consequently satisfies
\[
h^0\!\left(X_C,\OO_{X_C}(f_C^{-1}(p))\right)\ge1+g^g
\]
for every complex structure and every point.
\end{theorem}

The uniformization is a finite-type argument, but two details are essential.  Schneider's stability theorem must be applied to the canonical Fourier--Mukai Raynaud bundle together with its theta-group equivariance.  On a finite level cover of moduli, the mapping class subgroup must fix each actual homomorphism $\rho_i$, not only its conjugacy class; otherwise a projective descent cocycle remains.  With these points in place, there are only finitely many possible marked representations, their all-points failure loci are Zariski closed, and irreducibility of the level cover forces one representation to work everywhere.

For $g\ge3$, Theorem~\ref{thm:intro-uniform} gives a finite-unitary counterexample to Litt's generic-global-generation conjecture \cite[Conjecture~6.3.5]{Litt2024}: the bundle $V_{\rho_g,C}^{\vee}\otimes K_C$ is not generated at any point for any complex structure.  Since $|G_g|\ge1+g^{2g}>g^2$, the small-deck-group theorem of Landesman--Litt \cite[Theorem~1.2]{LandesmanLitt2023} does not decide the Putman--Wieland orbit question for this cover.  Failure of generic global generation alone does not produce a finite orbit.

The same evaluation map gives a numerical constraint that is independent of Raynaud's construction.

\begin{theorem}[Degree bounds]\label{thm:intro-degree}
Let $C$ have genus $g\ge2$.  If a nontrivial stable degree-zero bundle $W$ of rank $r$ satisfies $H^0(C,W(p))\ne0$ for every $p$, and
\[
c=\min_{p\in C}h^0(C,W(p)),
\]
then $r\ge gc+1$.  Consequently, every connected Prill-exceptional finite \'etale cover of $C$ has degree at least $g+2$; in the Galois case the degree is at least $1+(g+1)^2$.
\end{theorem}

The Putman--Wieland connection is global and weight one.  The existence of a nonzero finite orbit is equivalent to the existence, after finite \'etale base change, of a nonzero constant rational subvariation.  The following rank obstruction therefore separates special isoperiodic fibres from finite-index monodromy.

\begin{theorem}[Maximal-rank fixed-part obstruction]\label{thm:intro-fixed}
Let $S$ be a connected smooth complex quasi-projective variety, and let $\mathbb V_{\Z}$ be a polarized integral variation of Hodge structure of type $\{(1,0),(0,1)\}$ and rank $2n$.  If its local period map has differential of rank $n(n+1)/2$ at one point of $S$, then $\mathbb V_{\Q}$ has no nonzero constant subvariation after any connected finite \'etale base change.
\end{theorem}

For the Prym map $P_6:\cR_6\to\cA_5$, both spaces have dimension $15$ and the generic degree is $27$ \cite[Theorem~I.2.1]{DonagiSmith1981}.  Theorem~\ref{thm:intro-fixed} excludes a finite orbit in its anti-invariant cohomology and gives a no-extension statement for flat classes on the smooth-Hurwitz part of a special Fano fibre.  The same argument applies to the degree-$3$ map $P_{3,4}$ and its special Eckardt fibre.  These are independent period-rank proofs of double-cover cases already covered by \cite[Theorem~1.2]{LandesmanLitt2023}; the ramified case belongs to the punctured-surface formulation and is not an unramified cover of a closed genus-three surface.

The limitation is structural.  The $F_4$ calculation gives a nonisotrivial weight-two summand of rank $26$ over a cyclotomic field.  The uniform Raynaud representation gives failure of generic global generation in weight one.  A Putman--Wieland counterexample requires a rational flat class after finite-index base change.  Neither construction supplies that additional structure, while Theorem~\ref{thm:intro-fixed} rules it out for the maximal-rank Prym maps considered here.

\section{The Eckardt quotient and its \texorpdfstring{$F_4$}{F4} Tannaka group}\label{sec:f4}

All varieties in Sections~\ref{sec:f4}--\ref{sec:eckardt-period} are over $\C$.  Constructible complexes have complex coefficients.  For an integral subvariety $Z$ of an abelian variety, $\delta_Z$ denotes its perverse intersection complex, extended by zero.  The convolution Tannaka group $G(\delta_Z)$ is taken in the quotient by negligible complexes, following \cite{KramerWeissauer2015}.

Consider
\begin{equation}\label{eq:Y0}
Y_0=\{F_0=0\}\subset\PP^4,
\qquad
F_0=x_0^2x_1+x_1^3+x_2^3+x_3^3+x_4^3,
\end{equation}
and $p_0=[1:0:0:0:0]$.  The five partial derivatives are
\[
2x_0x_1,\quad x_0^2+3x_1^2,\quad3x_2^2,\quad3x_3^2,\quad3x_4^2.
\]
They have no common projective zero, so $Y_0$ is smooth.  The tangent hyperplane at $p_0$ is $x_1=0$, and its section is the cone over the smooth Fermat plane cubic $x_2^3+x_3^3+x_4^3=0$.  Thus $p_0$ is an Eckardt point.

Let $F=F(Y_0)$ be the Fano surface of lines.  Its Albanese variety $A$ has dimension five, and the Albanese map is an embedding \cite{ClemensGriffiths1972}.  The lines through $p_0$ form a smooth elliptic curve $E\subset F$.  Roulleau's reflection $\iota:F\to F$ fixes $E$ pointwise, acts by $+1$ on $T_0E$, and by $-1$ on $T_0A/T_0E$ \cite{Roulleau2009}.  Normalize the Albanese embedding by translating by a point of $E$ so that every coplanar triple of lines has sum zero.  This normalization leaves $E\subset A$ unchanged and gives
\begin{equation}\label{eq:reflection-quotient}
u\circ\iota=[-1]_B\circ u,
\qquad B=A/E,
\end{equation}
where $u:F\to B$ is the quotient map.

\begin{proposition}[The Eckardt $F_4$ quotient]\label{prop:f4-quotient}
Put $X=u(F)\subset B$.  Then $u:F\to X$ is birational, its only positive-dimensional fibre is $E=u^{-1}(0)$, and
\begin{equation}\label{eq:semismall}
Ru_*\C_F[2]\simeq\delta_X\oplus\delta_0,
\qquad \chi(B,\delta_X)=26.
\end{equation}
Moreover,
\begin{equation}\label{eq:f4-group}
G(\delta_X)\simeq F_4,
\qquad
\mathbf{27}|_{F_4}\simeq\mathbf1\oplus\mathbf{26}.
\end{equation}
Here $\mathbf{27}$ is the Albert representation of the simply connected group $E_6$ and $\mathbf{26}$ is the minimal representation of $F_4$.
\end{proposition}

\begin{proof}
The tangent-bundle theorem identifies $\PP(T_sF)$ with the line $\ell_s\subset Y_0$ represented by $s$, inside $\PP(T_0A)=\PP^4$ \cite{ClemensGriffiths1972}.  Since $\PP(T_0E)=\{p_0\}$ \cite[Proposition~4]{Roulleau2009}, the differential of $u$ has a kernel precisely at the points representing lines through $p_0$, namely on $E$.  Thus $E$ is the only contracted curve.

To compute the generic degree, project the lines $\ell_s$ from $p_0$.  If $u(s)=u(t)$ at a general point outside $B[2]$, the tangent-plane identity shows that $\ell_s$ and $\ell_t$ lie in the same plane through $p_0$.  The Eckardt plane-section decomposition says that the only possibilities away from $E$ are $t=s$ and $t=\iota(s)$.  In the second case \eqref{eq:reflection-quotient} gives $u(t)=-u(s)$, forcing $u(s)\in B[2]$.  Hence the generic degree is one.

The map is semismall.  Its unique relevant exceptional stratum is the origin, and $E^2=-3$ \cite[Proposition~10]{Roulleau2009}.  Thus the local intersection form is nondegenerate of rank one.  The semismall decomposition theorem \cite{deCataldoMigliorini2002} therefore gives exactly one skyscraper summand.  Since $\chi(F)=27$ \cite{ClemensGriffiths1972}, taking Euler characteristics proves \eqref{eq:semismall}.

For completeness, I recall the group-theoretic identification.  With the coplanar-triple normalization, the incidence correspondence supplies an invariant cubic in the convolution category.  Kr\"amer's Gauss-monodromy calculation identifies the connected derived Tannaka group of $\delta_F$ with $E_6$ on its minuscule $27$-dimensional module \cite{Kramer2016}; the invariant cubic and its full stabilizer then give $G(\delta_F)=E_6$ \cite{Kramer2022}.  Functoriality under $A\to B$ and \eqref{eq:semismall} give a closed subgroup $H=G(\delta_X)\subset E_6$ with
\[
\mathbf{27}|_H=\mathbf1\oplus V,
\qquad \dim V=26.
\]
The translation stabilizer of $X$ is trivial.  To see this, factor $u$ through the normalization $F\to X^\nu\to X$.  The first map is an isomorphism away from $E$ and has no other exceptional curve.  The image $z$ of $E$ is singular: over a smooth surface point, the exceptional curves of a birational morphism from a smooth surface arise from point blow-ups and are rational, whereas $E$ is elliptic.  Thus $z$ is the unique singular point of $X^\nu$.  A translation preserving $X$ lifts to $X^\nu$, fixes $z$, and therefore fixes its image $0\in B$; the translation is zero.  Thus $\delta_X$ is nondivisible in the sense of Javanpeykar--Kr\"amer--Lehn--Maculan.  Their irreducibility theorem \cite[Proposition~3.3(2)]{JKLM2025} implies that $V|_{H^\circ}$ is irreducible.

Let $v_0$ span the fixed line in the Albert algebra $J$.  If its cubic norm vanished, then $v_0$ would have Albert rank one or two.  The Hessian of the Albert cubic has kernel dimension respectively $17$ or $9$ at standard representatives of these two orbits.  On the other hand, $H^\circ$-equivariance and $J|_{H^\circ}=\mathbf1\oplus V$ allow only the kernel dimensions $0,1,26,27$, a contradiction.  Thus the norm of $v_0$ is nonzero.  Its stabilizer in $E_6$ is $F_4$ \cite{SpringerVeldkamp2000}, so $H\subset F_4$.  Finally, no proper positive-dimensional closed subgroup of $F_4$ in characteristic zero acts irreducibly on the minimal $26$-dimensional module \cite[Theorem~1 and Table~1.2]{LiebeckSeitz2004}.  Hence $H^\circ=F_4$, and $F_4=H^\circ\subset H\subset F_4$.  This proves \eqref{eq:f4-group}.
\end{proof}

\begin{remark}\label{rem:f4-not-monodromy}
Proposition~\ref{prop:f4-quotient} identifies a fibrewise convolution Tannaka group.  It does not identify the connected geometric monodromy of a variation with $F_4$.  Such an identification requires an additional comparison of Hodge tensors or an equivalent normality argument.
\end{remark}

\section{The marked Eckardt branch and the reflection bound}\label{sec:eckardt-algebra}

Let
\[
\iota(x_0,x_1,x_2,x_3,x_4)=(-x_0,x_1,x_2,x_3,x_4).
\]
The projective involution of $Y_0$ induces Roulleau's reflection on $F(Y_0)$ by the plane-section characterization, and both involutions are denoted by $\iota$.
At the highly symmetric cubic \eqref{eq:Y0}, the unmarked Eckardt locus has several local branches.  Throughout this section, $T_E$ means the tangent space at $(Y_0,p_0)$ to the local moduli space of \emph{marked} Eckardt pairs, equivalently the local fixed branch determined by $p_0$.

The Jacobian ring is
\begin{equation}\label{eq:jacobian-ring}
R=\C[x_0,\ldots,x_4]/(x_0x_1,x_0^2+3x_1^2,x_2^2,x_3^2,x_4^2).
\end{equation}
The standard Jacobian-ring identification $H^1(Y_0,T_{Y_0})\simeq R_3$, composed with the deformation isomorphism $H^1(Y_0,T_{Y_0})\simeq H^1(F,T_F)$, gives
$H^1(F,T_F)\simeq R_3$; see the proof of \cite[Proposition~3.8]{KLM2026}.  Define
\begin{align*}
a_0&=x_2x_3x_4,&
a_1&=x_1x_3x_4,&
a_2&=x_1x_2x_4,&
a_3&=x_1x_2x_3,\\
a_4&=x_1^2x_4,&
a_5&=x_1^2x_3,&
a_6&=x_1^2x_2.
\end{align*}

\begin{lemma}\label{lem:eckardt-tangent}
Under $H^1(F,T_F)\simeq R_3$,
\[
T_E=R_3^{\iota}=\langle a_0,\ldots,a_6\rangle.
\]
In particular, $\dim T_E=7$.
\end{lemma}

\begin{proof}
After shrinking, choose an $\iota$-equivariant analytic $\operatorname{PGL}_5$-slice: the uniqueness of the reflection attached to a marked Eckardt point permits every nearby reflection to be analytically conjugated to the fixed involution $\iota$, and the marked point to $p_0$.  Every marked Eckardt deformation in this slice has an equation of the form
\[
x_0^2\ell(x_1,\ldots,x_4)+C(x_1,\ldots,x_4),
\]
where $\ell\ne0$ at a smooth marked Eckardt pair.  Conversely, every smooth invariant deformation of this form has $p_0$ as a marked Eckardt point: its tangent section $\ell=0$ is the cone with vertex $p_0$ over $C|_{\ell=0}=0$.  Hence the deformation functor of the marked branch is the $\iota$-fixed subfunctor.  A finite-order fixed locus is smooth in characteristic zero, and its tangent space is $R_3^{\iota}$.  Reduction of cubic monomials by \eqref{eq:jacobian-ring} gives exactly the seven displayed basis elements.
\end{proof}

Let
\begin{equation}\label{eq:mu}
\mu:\Sym^2H^1(F,T_F)\longrightarrow H^2(F,\wedge^2T_F)
\end{equation}
be cup product.  Kr\"amer--Litt--Maculan identify the relevant coefficient of $\mu$ with
\begin{equation}\label{eq:nu}
\nu:\Sym^2R_3\longrightarrow
\Hom(\wedge^2R_1,\wedge^2R_4),
\qquad
\nu(f\odot g)(r\wedge s)=fr\wedge gs+gr\wedge fs.
\end{equation}
This is the Jacobian-ring formula in the proof of \cite[Proposition~3.8]{KLM2026}.

\begin{proposition}\label{prop:rank27}
The restricted cup product has kernel
\begin{equation}\label{eq:kernel-mu}
\ker(\mu|_{\Sym^2T_E})
=\left\langle a_1a_6+a_2a_5+a_3a_4\right\rangle.
\end{equation}
Consequently,
\[
I_E:=\mu(\Sym^2T_E)
\quad\text{has dimension }27.
\]
\end{proposition}

\begin{proof}
In the bases of Appendix~\ref{app:certificates}, exact row reduction of \eqref{eq:nu} gives the relation in \eqref{eq:kernel-mu}.  The appendix also records a $27\times27$ minor with determinant $-128$, so this is the entire kernel of $\nu|_{\Sym^2T_E}$.

It remains to pass from $\nu$ to $\mu$.  The binary cubic $x_0^2x_1+x_1^3$ has three distinct roots, so $Y_0$ is projectively equivalent to the Fermat cubic.  At the Fermat cubic, \eqref{eq:nu} has rank $50$ \cite[Proposition~3.6]{KLM2026}.  The bicanonical section of lines of second type lies in $\ker\mu^\vee$, so $\rk\mu\le50$ \cite[Proposition~3.7]{KLM2026}.  Since $\nu$ factors through $\mu$, both maps have rank $50$ and therefore have the same kernel.  Restriction to $\Sym^2T_E$ proves the proposition.
\end{proof}

Let $\widehat B=\Pic^0(B)$ and take $0\ne M\in\widehat B$.  Let $\mathbb M$ be the unique unitary rank-one local system whose holomorphic bundle is $M$, and put
\[
\mathbb L=u^*\mathbb M,
\qquad L=u^*M.
\]
For $0\le p\le2$, write
\[
H_L^p=H^{2-p}(F,\Omega_F^p\otimes L).
\]
Put $H_L^p=0$ otherwise.  The extreme Hodge spaces satisfy
\[
H_L^2=E_L=H^0(F,K_F\otimes L),
\qquad
(H_L^0)^\vee=E_{L^{-1}}=H^0(F,K_F\otimes L^{-1}).
\]
They both have dimension six by \cite[Lemma~3.3]{KLM2026}.  Let
\[
\pi_E:H^0(F,K_F^2)\longrightarrow I_E^\vee
\]
be the Serre-dual projection.

On a simply connected analytic neighborhood of $(Y_0,p_0)$ in the marked Eckardt branch, use an Ehresmann trivialization of the pair $(F,E)$ to transport the character of $\mathbb L$.  It remains trivial on $E_s$ and therefore is pulled back from $B_s=\Alb(F_s)/E_s$.  The relative twisted cohomology is then a complex variation of Hodge structure.  For $\xi\in T_E$, its infinitesimal Higgs field
\[
\theta_L(\xi):H_L^p\longrightarrow H_L^{p-1}
\]
is contraction and cup product with the Kodaira--Spencer class $\xi$.  Its top-to-bottom second iterate factors through \eqref{eq:mu}; under Serre duality its coefficient map is the projected multiplication used below, as in \cite[Equation~(2.1) and the proof of Proposition~2.6]{KLM2026}.

\begin{proposition}[Reflection bound]\label{prop:reflection-bound}
For every $0\ne M\in\widehat B$, the top-to-bottom second Higgs iterate along $T_E$ satisfies
\begin{equation}\label{eq:upper21}
\rk\theta_L^{(2)}|_{\Sym^2T_E}\le21.
\end{equation}
\end{proposition}

\begin{proof}
Equation \eqref{eq:reflection-quotient} gives $\iota^*L\simeq L^{-1}$.  Rigidify $M$ at the origin.  The theorem of the cube gives a rigidified isomorphism $[-1]^*M\simeq M^{-1}$; after pullback it may be normalized so that applying $\iota$ twice is the identity.  Because $\iota$ fixes $T_E$ pointwise, equivariance of cup product shows that it fixes $I_E$ pointwise.  Hence $\pi_E\iota^*=\pi_E$.

Choose the involutive normalization
\[
\rho_L=\iota^*:E_L\xrightarrow{\sim}E_{L^{-1}},
\qquad \rho_{L^{-1}}\rho_L=1,
\]
and define
\[
b_L(s,t)=\pi_E\bigl(s\,\rho_L(t)\bigr).
\]
Then
\[
b_L(t,s)=\pi_E\!\left(\iota^*(t\rho_L(s))\right)
=\pi_E(\rho_L(t)s)=b_L(s,t),
\]
so $b_L$ factors through $\Sym^2E_L$.  By \cite[Equation~(2.1) and the proof of Proposition~2.6]{KLM2026}, the transpose of the coefficient map of $\theta_L^{(2)}$ on $I_E$ is exactly the projected multiplication $E_L\otimes E_{L^{-1}}\to I_E^\vee$.  Therefore
\[
\rk\theta_L^{(2)}|_{\Sym^2T_E}
=\rk b_L
\le\dim\Sym^2E_L
=21.
\]
\end{proof}

\section{Maximal infinitesimal variation on the marked branch}\label{sec:eckardt-period}

The residue action of $\iota$ includes the determinant character: under
$H^0(F,\Omega_F^1)\simeq R_1$, the vector $x_0$ is invariant and
$x_1,\ldots,x_4$ are anti-invariant.  Let
\begin{equation}\label{eq:alpha}
\alpha\in R_1^\vee,
\qquad
\alpha(x_0)=0,
\qquad
\alpha(x_1)=\alpha(x_2)=\alpha(x_3)=\alpha(x_4)=1.
\end{equation}
The annihilation of the invariant line means that
$\alpha\in T_0\widehat B\subset T_0\Pic^0(F)$.

Choose a smooth analytic arc $\Delta\to\widehat B$ through zero with tangent
$\alpha$, and let $\mathscr L$ be the induced Poincar\'e line bundle on
$F\times\Delta$.  Write $p:F\times\Delta\to\Delta$ and set
\begin{equation}\label{eq:direct-images}
\mathscr E=p_*(K_F\otimes\mathscr L),
\qquad
\mathscr E^-=p_*(K_F\otimes\mathscr L^{-1}).
\end{equation}

\begin{lemma}[Saturated central fibre]\label{lem:saturated-limit}
The sheaves $\mathscr E$ and $\mathscr E^-$ are locally free of rank six.  Their central fibres inject into $H^0(F,K_F)$ and satisfy
\begin{equation}\label{eq:Ualpha}
\mathscr E_0=\mathscr E^-_0
=U_\alpha:=\wedge^2\ker\alpha,
\qquad \dim U_\alpha=6.
\end{equation}
Multiplication of relative sections defines a regular morphism
\begin{equation}\label{eq:relative-multiplication}
\mathscr E\otimes\mathscr E^-
\longrightarrow H^0(F,K_F^2)\otimes\OO_\Delta.
\end{equation}
In particular, its central value requires neither a pole nor a renormalization.
\end{lemma}

\begin{proof}
Both direct images are torsion-free on the smooth curve $\Delta$, hence locally free.  Away from zero they have rank six by \cite[Lemma~3.3]{KLM2026}.  A relative section whose restriction at zero vanishes is divisible by a local parameter, so restriction embeds each central fibre into $H^0(F,K_F)$.  The first obstruction to extending a canonical form is cup product with $\alpha$ for $\mathscr E$ and with $-\alpha$ for $\mathscr E^-$.  Under the residue identifications, \cite[Lemma~3.4]{KLM2026} identifies this kernel with $\wedge^2\ker\alpha$, which has dimension six.  Rank and injectivity give \eqref{eq:Ualpha}.  Ordinary multiplication on $F\times\Delta$ gives \eqref{eq:relative-multiplication}; all sheaves involved are locally free, so the morphism is regular at zero.
\end{proof}

Compose \eqref{eq:relative-multiplication} with $\pi_E$, transpose, and precompose with $\mu|_{\Sym^2T_E}$.  This gives a morphism of vector bundles
\begin{equation}\label{eq:relative-second}
\overline\Theta^{(2)}:
\Sym^2T_E\otimes\OO_\Delta
\longrightarrow
\mathcal Hom(\mathscr E,(\mathscr E^-)^\vee).
\end{equation}
For $t\ne0$, Serre duality and \cite[Equation~(2.1)]{KLM2026} identify it with the top-to-bottom second Higgs iterate.  At zero, it is the saturated limit obtained by restricting ordinary multiplication to
$U_\alpha\otimes U_\alpha$.  Notice that
$\overline\Theta^{(2)}_0$ is not the second Higgs field of the untwisted ten-dimensional space $H^0(F,K_F)$.

For the exact computation, put
\[
z=x_0,
\qquad
w_1=x_1-x_4,
\quad w_2=x_2-x_4,
\quad w_3=x_3-x_4,
\]
and order a basis of $U_\alpha$ as
\begin{equation}\label{eq:ubasis}
\begin{array}{lll}
u_0=z\wedge w_1,&u_1=z\wedge w_2,&u_2=z\wedge w_3,\\
u_3=w_1\wedge w_2,&u_4=w_1\wedge w_3,&u_5=w_2\wedge w_3.
\end{array}
\end{equation}
Formula \eqref{eq:nu} computes the matrix of \eqref{eq:relative-second} at zero.  The complete certificates are recorded in Appendix~\ref{app:certificates}.

\begin{lemma}[The two certificate minors]\label{lem:minors}
The map $\overline\Theta^{(2)}_0$ has rank exactly $12$.  Moreover, the flattening
\begin{equation}\label{eq:flattening}
\Phi_0:T_E\longrightarrow
\Hom\!\left(T_E,\Hom(U_\alpha,U_\alpha^\vee)\right),
\qquad
\Phi_0(\xi)(\eta)=\overline\Theta^{(2)}_0(\xi\eta),
\end{equation}
is injective.
\end{lemma}

\begin{proof}
The involution splits $U_\alpha=U_\alpha^+\oplus U_\alpha^-$ with dimensions $3+3$.  Since $I_E$ is fixed, the mixed part of the projected multiplication is zero.  At $t=0$ the two direct-image fibres coincide with $U_\alpha$, and ordinary multiplication is commutative; hence the surviving same-parity blocks factor through $\Sym^2U_\alpha^+\oplus\Sym^2U_\alpha^-$.  Its rank is therefore at most
$\dim\Sym^2U_\alpha^++\dim\Sym^2U_\alpha^-=6+6=12$.
The matrix in \eqref{eq:matrix12} is a $12\times12$ minor and has determinant
$-2^7 3^6$, proving equality.

For \eqref{eq:flattening}, the seven rows and columns displayed in
\eqref{eq:matrix7} give a determinant $-2^3 3^7$.  Thus $\Phi_0$ has rank seven and is injective.
\end{proof}

\begin{theorem}[Maximal variation on the marked Eckardt branch]\label{thm:eckardt-main}
There is a nonempty Zariski-open subset
\[
\Omega\subset\widehat B\setminus\{0\}
\]
such that, for every $M\in\Omega$ and $L=u^*M$, the full infinitesimal Higgs map
\begin{equation}\label{eq:first-higgs}
\theta_L|_{T_E}:T_E\longrightarrow
\bigoplus_p\Hom(H_L^p,H_L^{p-1})
\end{equation}
is injective, and
\begin{equation}\label{eq:rank-range}
12\le\rk\theta_L^{(2)}|_{\Sym^2T_E}\le21.
\end{equation}
The open set $\Omega$ contains torsion points of arbitrarily large order.
\end{theorem}

\begin{proof}
Extend the flattening in \eqref{eq:flattening} to $\Phi_t$ by using
\eqref{eq:relative-second}.  Lemma~\ref{lem:minors} and openness of matrix rank imply that $\Phi_t$ is injective and that the second iterate has rank at least $12$ for all sufficiently small $t\ne0$.  Higgs integrability gives
\[
\theta_{L_t}^{(2)}(\xi,\eta)
=\theta_{L_t}(\eta)\theta_{L_t}(\xi)
=\theta_{L_t}(\xi)\theta_{L_t}(\eta).
\]
Therefore $\ker\theta_{L_t}\subseteq\ker\Phi_t$, and the first Higgs field is injective along $T_E$.  Proposition~\ref{prop:reflection-bound} gives the upper bound $21$.

On $\widehat B\setminus\{0\}$, \cite[Lemma~3.3]{KLM2026} and cohomology and base change make the two Poincar\'e direct images algebraic vector bundles of rank six.  Multiplication and flattening are algebraic morphisms.  The two rank conditions define a Zariski-open locus; the punctured analytic arc meets it by Lemma~\ref{lem:minors}, so it is nonempty.  Denote this locus by $\Omega$.  Torsion points are Zariski dense in an abelian variety.  Since the torsion points of bounded order form a finite set, every nonempty open subset contains torsion points of arbitrarily large order.
\end{proof}

\begin{corollary}[The $26$-dimensional torsion summand]\label{cor:26summand}
Let $M\in\Omega$ have finite order $n$, let $K=\Q(\zeta_n)$, and let
$\chi:\pi_1(B)\to\mu_n\subset K^\times$ be the character determined by $M$.  Let $\tau_0:K\hookrightarrow\C$ be the embedding for which the associated holomorphic line bundle is $M$.  After passage to a finite level cover and restriction to a simply connected analytic neighborhood $S_E$ in the local marked Eckardt moduli space, write $\pi:\mathscr F\to S_E$ for the resulting Fano-surface family and $\mathbb L_\chi$ for the transported pullback of $\chi$, rigidified to be trivial on the relative elliptic divisor.  Then the $K$-local system
\begin{equation}\label{eq:tate-splitting}
\mathbb V_\chi=R^2\pi_*\mathbb L_\chi
\simeq K(-1)\oplus\mathbb W_\chi,
\qquad \rk_K\mathbb W_\chi=26
\end{equation}
as a horizontal direct sum.  For every embedding $\tau:K\hookrightarrow\C$, the Hodge numbers of
$\mathbb W_{\chi,\tau}=\mathbb W_\chi\otimes_{K,\tau}\C$ are $(6,14,6)$.  For $\tau=\tau_0$, its period differential along $T_E$ has rank seven.  Consequently,
$\operatorname{Res}_{K/\Q}\mathbb W_\chi$ has $\Q$-rank $26[K:\Q]$ and period differential of rank seven along $T_E$.
\end{corollary}

\begin{proof}
Let $i:\mathscr E\hookrightarrow\mathscr F$ be the relative elliptic divisor.  Fibrewise triviality shows that $\mathbb L_\chi|_{\mathscr E}$ is pulled back from $S_E$; since $S_E$ is simply connected, the chosen rigidification makes this restriction trivial.  Relative Gysin and restriction are morphisms of variations of Hodge structure
\[
K(-1)\xrightarrow{i_*}R^2\pi_*\mathbb L_\chi
\xrightarrow{i^*}K(-1).
\]
Their composite is multiplication by the fibrewise self-intersection $E_s^2=-3$.  Hence
\[
P=-\frac13\,i_*\circ i^*
\]
is a horizontal idempotent.  Its image is $K(-1)$ and its kernel is the rank-$26$ variation $\mathbb W_\chi$, proving \eqref{eq:tate-splitting}.

For every embedding $\tau$, the character $\tau\chi$ is nontrivial.  The proof of \cite[Lemma~3.3]{KLM2026} gives
\[
H^j(F,\mathbb L_{\chi,\tau})=0\qquad(j\ne2).
\]
Thus $\dim_\C H^2(F,\mathbb L_{\chi,\tau})=\chi(F)=27$.  Its two extreme Hodge numbers are six by the same lemma, so its Hodge numbers are $(6,15,6)$.  Removing the Tate line gives $(6,14,6)$ for every $\tau$.

For the distinguished embedding $\tau_0$, Theorem~\ref{thm:eckardt-main} gives rank seven; the Tate line has zero Higgs field, so this rank occurs on $\mathbb W_{\chi,\tau_0}$.  After restriction of scalars the period differential contains this component and has source $T_E$ of dimension seven, hence has rank exactly seven.
\end{proof}

\begin{remark}\label{rem:coeff-field}
The rank $26$ in Corollary~\ref{cor:26summand} is over the character field $K$, or over $\C$ after choosing one embedding.  The underlying rational variation has rank $26[K:\Q]$, not $26$ in general.  Moreover, the result is a weight-two statement.  These two facts are essential when comparing it with Putman--Wieland.
\end{remark}

\section{Raynaud bundles and finite monodromy}\label{sec:raynaud}

Let $C$ be a smooth projective connected curve of genus $g\ge2$, let
$J=\Pic^0(C)$, and fix a symmetric theta divisor $\Theta$.  Let
$i:C\hookrightarrow J$ be an Abel--Jacobi embedding.  With the conventional harmless inversion depending on the Fourier--Mukai normalization, define the canonical Raynaud bundle
\begin{equation}\label{eq:FM-raynaud}
\mathscr E_n=\mathcal F_J\bigl(\OO_J(-n\Theta)\bigr)[g].
\end{equation}
Set $R_n=i^*\mathscr E_n$.  Raynaud proves that $R_n$ is semistable and
\begin{equation}\label{eq:kummer-raynaud}
[n]^*\mathscr E_n\simeq
\OO_J(n\Theta)\otimes H^0(J,\OO_J(n\Theta))^\vee,
\end{equation}
and
\begin{equation}\label{eq:raynaud-properties}
\rk R_n=n^g,
\qquad \mu(R_n)=\frac gn,
\qquad H^0(C,R_n\otimes\alpha)\ne0
\quad(\alpha\in\Pic^0(C)).
\end{equation}
The last assertion follows on all of $\Pic^0(C)$ from Raynaud's generic assertion by upper semicontinuity \cite{Raynaud1982}.

\begin{lemma}[Stability]\label{lem:raynaud-stable}
The bundle $R_n=i^*\mathscr E_n$ is stable.
\end{lemma}

\begin{proof}
This is \cite[Corollary~2.2]{Schneider2007}; I recall why the canonical Fourier--Mukai structure matters.  Formula \eqref{eq:kummer-raynaud} is equivariant for Mumford's theta group of $\OO_J(n\Theta)$.  Its weight-one Schr\"odinger module $H^0(J,\OO_J(n\Theta))$ is irreducible, and hence so is its dual \cite[Proposition~2.1(3)]{Schneider2007}.

The pullback
\[
\widetilde C=C\times_{J,[n]}J
\]
is connected because $i_*:\pi_1(C)\to\pi_1(J)=H_1(C,\Z)$ is abelianization.  If a saturated proper subbundle of $R_n$ had the same slope, its pullback, after tensoring by the inverse of the line factor in \eqref{eq:kummer-raynaud}, would be a degree-zero subbundle of a trivial bundle on $\widetilde C$.  Such a subbundle is constant: dualize, choose a full set of generating sections, and use degree zero to see that their determinant never vanishes.  Descent would therefore give a proper invariant subspace of the Schr\"odinger module, contradicting irreducibility.  Raynaud's semistability then implies stability.
\end{proof}

Set
\[
r=g^g,
\qquad R=R_g.
\]
Choose $M\in\Pic^{-1}(C)$ such that
\begin{equation}\label{eq:det-root}
M^{\otimes r}\simeq(\det R)^{-1};
\end{equation}
it exists because multiplication by $r$ on $J$ is surjective.  Put
\begin{equation}\label{eq:VC}
V_C=R\otimes M.
\end{equation}
Then $V_C$ is stable of degree zero with trivial determinant.  A nonzero section would produce a line subbundle of nonnegative degree, contrary to stability, so
\begin{equation}\label{eq:VC-properties}
H^0(C,V_C)=0,
\qquad
H^0(C,V_C(p))\ne0\quad(p\in C),
\end{equation}
where the second assertion follows from \eqref{eq:raynaud-properties} because $M(p)\in\Pic^0(C)$.

\begin{proposition}[A bounded trivializing cover]\label{prop:trivializing-cover}
The bundle $V_C$ is trivialized by a connected finite \'etale cover of degree at most $g^{3g}$.  Its Narasimhan--Seshadri representation is irreducible, unitary, has trivial determinant, and has finite image of order at most $g^{3g}$.
\end{proposition}

\begin{proof}
Pull $C\to J$ back along $[g]:J\to J$.  The resulting connected Galois cover
\[
q:\widetilde C\longrightarrow C
\]
has degree $g^{2g}$: its subgroup is the kernel of
$\pi_1(C)\to H_1(C,\Z/g)$.  Equations \eqref{eq:kummer-raynaud} and \eqref{eq:VC} give
\[
q^*V_C\simeq N^{\oplus r},
\qquad
N=\widetilde i^*\OO_J(g\Theta)\otimes q^*M.
\]
Taking determinants and using $\det V_C\simeq\OO_C$ gives $N^{\otimes r}\simeq\OO_{\widetilde C}$.  If $d\mid r$ is the exact order of $N$, the algebra
\[
\bigoplus_{a=0}^{d-1}N^{-a}
\]
defines an \'etale $\mu_d$-torsor trivializing $N$.  For $0<a<d$, the degree-zero line bundle $N^a$ is nontrivial and has no section.  Therefore
\[
H^0\!\left(\widetilde C,\bigoplus_{a=0}^{d-1}N^{-a}\right)=\C.
\]
Its coordinate algebra has no nontrivial idempotent, so the torsor is connected.  The composite cover has degree
\[
d g^{2g}\le g^g g^{2g}=g^{3g}.
\]

By the Narasimhan--Seshadri correspondence \cite{NarasimhanSeshadri1965}, the stable degree-zero bundle $V_C$ corresponds to an irreducible unitary representation $\rho_C$.  Pullback to the displayed cover is the trivial polystable bundle, hence the restriction of $\rho_C$ to the covering subgroup is trivial.  Therefore its image is finite and its order is at most the degree of the cover.  The determinant character of $\rho_C$ is the rank-one Narasimhan--Seshadri representation corresponding to $\det V_C\simeq\OO_C$, hence is trivial.  Thus $\rho_C$ takes values in $\SU(r)$.
\end{proof}


\section{A single Prill representation for every complex structure}\label{sec:uniform}

Fix a pointed oriented closed surface $(\Sigma_g,x_0)$ and write
$\pi=\pi_1(\Sigma_g,x_0)$.  Proposition~\ref{prop:trivializing-cover} produces, for each marked complex structure, an irreducible finite representation of rank $r=g^g$ and image order at most $g^{3g}$.  I now show that one representation works for all complex structures.

\begin{lemma}[Finiteness of the possible representations]\label{lem:finite-reps}
Up to $\GL_r(\C)$-conjugacy, there are only finitely many nontrivial irreducible representations
\[
\rho:\pi\longrightarrow\SL_r(\C)
\]
whose image has order at most $g^{3g}$.
\end{lemma}

\begin{proof}
There are finitely many groups of bounded order: on a set of cardinality $m$, there are only finitely many multiplication tables.  A homomorphism from the $2g$-generated group $\pi$ to a fixed finite group $G$ is determined by the images of the generators, so there are at most $|G|^{2g}$ such homomorphisms.  Finally, a finite group has finitely many irreducible complex representations.  Retaining only faithful representations of dimension $r$ with trivial determinant proves the assertion.
\end{proof}

Average a positive Hermitian form over each finite image and choose representatives
\begin{equation}\label{eq:finite-list}
\rho_1,\ldots,\rho_N:\pi\longrightarrow\SU(r)
\end{equation}
of these conjugacy classes.  Let $\Gamma_{g,1}$ be the pointed mapping class group, acting on $\pi$ by actual automorphisms.  Since each $\rho_i$ takes values in a fixed finite matrix group, its orbit under precomposition is finite.  Define the exact stabilizer
\[
\Gamma_i=
\{\gamma\in\Gamma_{g,1}:\rho_i\circ\gamma_*=\rho_i\}.
\]
Choose $\ell\ge3$ and put
\begin{equation}\label{eq:exact-level}
\Gamma'=\Gamma(\ell)\cap\bigcap_{i=1}^N\Gamma_i.
\end{equation}
Then $\Gamma'$ is torsion-free and of finite index.  The quotient of pointed Teichm\"uller space
\[
S=T_{g,1}/\Gamma'
\]
is a connected smooth quasi-projective variety, finite \'etale over a fine level cover of $\cM_{g,1}$.  It carries a universal pointed curve
\begin{equation}\label{eq:universal-curve}
\pi_{\mathcal C}:\mathcal C\longrightarrow S.
\end{equation}

The section splits the homotopy sequence and gives
\[
\pi_1(\mathcal C)\simeq\pi\rtimes\Gamma'.
\]
Because \eqref{eq:exact-level} fixes $\rho_i$ as an actual homomorphism, not merely up to conjugacy,
\begin{equation}\label{eq:actual-descent}
\widetilde\rho_i(\delta,\gamma)=\rho_i(\delta)
\qquad(\delta\in\pi,\ \gamma\in\Gamma')
\end{equation}
is a homomorphism.  It defines a finite local system.  By the Riemann existence theorem the corresponding finite topological cover is algebraic, and finite-\'etale descent of its trivial bundle gives an algebraic flat bundle $\mathcal V_i$ on $\mathcal C$.  This exact descent is the reason for using the pointed mapping class group and the exact stabilizers in \eqref{eq:exact-level}.

Put
\[
\mathcal F_i=\mathcal V_i^\vee\otimes\omega_{\mathcal C/S}.
\]
Every fibre $V_{i,s}$ is a nontrivial irreducible unitary bundle of degree zero, so $H^0(V_{i,s})=0$.  Serre duality and Riemann--Roch give
\[
H^1(F_{i,s})=0,
\qquad h^0(F_{i,s})=r(g-1).
\]
Cohomology and base change make
\[
\mathcal H_i=(\pi_{\mathcal C})_*\mathcal F_i
\]
a vector bundle, and there is an evaluation map
\begin{equation}\label{eq:universal-evaluation}
\varepsilon_i:\pi_{\mathcal C}^*\mathcal H_i\longrightarrow\mathcal F_i.
\end{equation}
Let $D_i\subset\mathcal C$ be the determinantal closed subscheme defined by the maximal minors of $\varepsilon_i$; its support is the nonsurjectivity locus.

\begin{lemma}[The all-points locus is closed]\label{lem:all-points-closed}
The subset
\[
S_i=\{s\in S:C_s\subseteq D_i\}
\]
is Zariski closed.  A point $s$ lies in $S_i$ precisely when
\[
H^0(C_s,V_{i,s}(p))\ne0
\qquad\text{for every }p\in C_s.
\]
\end{lemma}

\begin{proof}
For $p\in C_s$, the cohomology sequence of
$0\to F_{i,s}(-p)\to F_{i,s}\to(F_{i,s})_p\to0$, followed by Serre duality, gives
\begin{equation}\label{eq:evaluation-duality}
\operatorname{coker}
\bigl[H^0(C_s,F_{i,s})\to(F_{i,s})_p\bigr]^\vee
\simeq H^0(C_s,V_{i,s}(p)).
\end{equation}
Thus the asserted pointwise equivalence holds.  The morphism $D_i\to S$ is proper.  Upper semicontinuity of fibre dimension shows that
$\{s:\dim(D_i)_s\ge1\}$ is closed.  Since $C_s$ is an irreducible curve, a closed subset of $C_s$ has dimension one exactly when its support is all of $C_s$.  This closed set is $S_i$.
\end{proof}

\begin{theorem}[Uniform Raynaud--Prill theorem]\label{thm:uniform}
Let $g\ge2$ and $r=g^g$.  There is a nontrivial irreducible finite representation
\begin{equation}\label{eq:rho-g}
\rho_g:\pi_1(\Sigma_g,x_0)\longrightarrow\SU(r)
\end{equation}
such that, for every marked compact Riemann surface $C$ of genus $g$ and every $p\in C$,
\begin{equation}\label{eq:uniform-property}
H^0(C,V_{\rho_g,C})=0,
\qquad
H^0(C,V_{\rho_g,C}(p))\ne0.
\end{equation}
If $G_g=\im\rho_g$, then
\begin{equation}\label{eq:group-bounds}
1+g^{2g}\le|G_g|\le g^{3g}.
\end{equation}
\end{theorem}

\begin{proof}
For each $s\in S$, the construction in Section~\ref{sec:raynaud} and Proposition~\ref{prop:trivializing-cover} supply an index $i$ for which $V_{i,s}$ is isomorphic to $V_{C_s}$ in \eqref{eq:VC}.  Hence \eqref{eq:VC-properties} and Lemma~\ref{lem:all-points-closed} give
\[
S=\bigcup_{i=1}^N S_i.
\]
The smooth connected variety $S$ is irreducible.  A finite union of proper closed subsets cannot equal an irreducible variety, so $S_j=S$ for some $j$.  Pulling back to Teichm\"uller space shows that the single marked representation $\rho_g=\rho_j$ has \eqref{eq:uniform-property} for every marked complex structure.

Every representation in the finite list has image order at most $g^{3g}$, proving the upper bound in \eqref{eq:group-bounds}; the determinant-one condition was built into the list.  By definition $G_g$ acts faithfully and irreducibly in dimension $r$.  Since $r>1$, the sum-of-squares formula for irreducible representations of a finite group contains both this representation and the distinct trivial representation, so
\[
|G_g|=\sum_{\tau\in\operatorname{Irr}(G_g)}(\dim\tau)^2
\ge1+r^2=1+g^{2g}.
\]
\end{proof}

\begin{corollary}[A fixed topological Prill cover]\label{cor:fixed-prill}
Let $f_C:X_C\to C$ be the connected Galois cover associated with $\ker\rho_g$.  It has the same topological covering type for every marked complex structure, degree $|G_g|$, and
\begin{equation}\label{eq:prill-many-sections}
h^0\!\left(X_C,\OO_{X_C}(f_C^{-1}(p))\right)
\ge1+g^g
\qquad(p\in C).
\end{equation}
\end{corollary}

\begin{proof}
The regular $G_g$-module contains the trivial representation once and the irreducible representation \eqref{eq:rho-g} with multiplicity $r$.  Therefore
\[
(f_C)_*\OO_{X_C}
\supset\OO_C\oplus V_{\rho_g,C}^{\oplus r}.
\]
Projection formula and \eqref{eq:uniform-property} give
\[
h^0\!\left(X_C,\OO_{X_C}(f_C^{-1}(p))\right)
\ge h^0(C,\OO_C(p))+r h^0(C,V_{\rho_g,C}(p))
\ge1+r.
\]
\end{proof}

\begin{corollary}[Failure of generic global generation]\label{cor:ggg}
For every $g\ge3$, the finite-unitary representation $\rho_g^\vee$ contradicts \cite[Conjecture~6.3.5]{Litt2024}.  More precisely, for every marked complex structure $C$, the bundle
\[
V_{\rho_g,C}^\vee\otimes K_C
\]
is not generated by global sections at any point of $C$.
\end{corollary}

\begin{proof}
Equation \eqref{eq:evaluation-duality}, with $V=V_{\rho_g,C}$, identifies the dual of the evaluation cokernel at $p$ with $H^0(C,V(p))$.  This space is nonzero at every point by Theorem~\ref{thm:uniform}.  The ggg conjecture asserts generic global generation for the holomorphic bundle associated with any fixed unitary representation on a generic complex structure.  The dual representation fails that assertion for every complex structure.
\end{proof}

\section{Numerical constraints on Prill-exceptional covers}\label{sec:degree}

The next estimate uses only stability, evaluation, and the vector-bundle Clifford inequality.  It applies well beyond the Raynaud construction.

\begin{theorem}[Rank--corank inequality]\label{thm:rank-corank}
Let $W$ be a nontrivial stable bundle of degree zero and rank $r$ on a curve $C$ of genus $g\ge2$.  Suppose
\[
H^0(C,W(p))\ne0\qquad(p\in C),
\]
and put
\[
c=\min_{p\in C}h^0(C,W(p)).
\]
Then
\begin{equation}\label{eq:rank-corank}
r\ge gc+1.
\end{equation}
\end{theorem}

\begin{proof}
Set $F=K_C\otimes W^\vee$.  Then $F$ is stable of rank $r$, degree $2r(g-1)$, and
\[
h^0(F)=r(g-1),
\]
because $H^1(F)^\vee=H^0(W)=0$.  Let $U\subset F$ be the saturation of the image of the evaluation map
\[
H^0(F)\otimes\OO_C\longrightarrow F,
\]
and write $D=\deg U$.  By upper semicontinuity, $h^0(C,W(p))$ takes its minimum value $c$ on a nonempty open subset of $C$.  At a point of this subset, \eqref{eq:evaluation-duality} says that the evaluation map to the $r$-dimensional fibre $F_p$ has cokernel of dimension $c$.  Its generic rank is therefore $r-c$, and saturation does not change generic rank; hence
\[
\rk U=r-c.
\]
All sections of $F$ factor through $U$, and $U\subset F$, so
\begin{equation}\label{eq:h0U}
h^0(U)=h^0(F)=r(g-1).
\end{equation}

The bundle $U$ is generically globally generated.  Every quotient bundle of $U$ is therefore generically globally generated, and the determinant of such a quotient has a nonzero section; hence every Harder--Narasimhan slope of $U$ is nonnegative.  The maximal-slope Harder--Narasimhan subbundle of $U$ is a proper subbundle of the stable bundle $F$, so its slope, and hence every Harder--Narasimhan slope of $U$, is strictly less than $2g-2$.  Apply the vector-bundle Clifford inequality
\[
h^0(Q)\le\frac{\deg Q}{2}+\rk Q
\]
to each semistable Harder--Narasimhan factor and add; see \cite[Theorem~2.1]{BPGN1997}.  Equation \eqref{eq:h0U} yields
\begin{equation}\label{eq:Dlower}
D\ge2r(g-2)+2c.
\end{equation}
On the other hand, stability of $F$ and integrality of degree give the strict bound
\begin{equation}\label{eq:Dupper}
D\le2(r-c)(g-1)-1.
\end{equation}
Combining \eqref{eq:Dlower} and \eqref{eq:Dupper} gives
$2cg\le2r-1$.  Thus $r>gc$, which is \eqref{eq:rank-corank}.
\end{proof}

\begin{corollary}[Degree bounds for covers]\label{cor:degree-bounds}
Let $f:X\to C$ be a connected finite \'etale cover of a genus-$g$ curve, $g\ge2$, and assume
\[
h^0\!\left(X,\OO_X(f^{-1}(p))\right)\ge2
\qquad(p\in C).
\]
Then
\begin{equation}\label{eq:nongalois-degree}
\deg f\ge g+2.
\end{equation}
If $f$ is Galois, then
\begin{equation}\label{eq:galois-degree}
\deg f\ge1+(g+1)^2.
\end{equation}
\end{corollary}

\begin{proof}
By projection formula,
\[
h^0\!\left(X,\OO_X(f^{-1}(p))\right)
=h^0\!\left(C,f_*\OO_X\otimes\OO_C(p)\right).
\]
The finite permutation local system decomposes as
\[
f_*\OO_X\simeq\OO_C\oplus\bigoplus_j W_j^{\oplus m_j},
\]
where the $W_j$ are nontrivial stable finite-unitary bundles.  Since $g\ge2$, one has $h^0(\OO_C(p))=1$.  Upper semicontinuity makes each locus
\[
Z_j=\{p\in C:H^0(C,W_j(p))\ne0\}
\]
closed.  The hypothesis and the displayed decomposition show that these finitely many loci cover $C$.  Irreducibility of $C$ gives $Z_j=C$ for some $j$.  Theorem~\ref{thm:rank-corank}, with $c\ge1$, gives $\rk W_j\ge g+1$, and
\[
\deg f=\rk(f_*\OO_X)\ge1+\rk W_j\ge g+2.
\]
If $f$ is Galois, set $r_j=\rk W_j$.  Its regular representation contains $W_j$ with multiplicity $r_j$.  Hence
\[
\deg f\ge1+r_j^2\ge1+(g+1)^2,
\]
because $r_j\ge g+1$.
\end{proof}

The same evaluation quotient measures a common infinitesimal Higgs kernel.  This gives a precise bridge to the local linear algebra used in Putman--Wieland arguments.

\begin{proposition}[Common Higgs kernel]\label{prop:common-higgs}
In the notation of Theorem~\ref{thm:rank-corank}, let $Q=F/U$ and define
\[
K_W^{1,0}
=\bigcap_{\xi\in H^1(C,T_C)}
\ker\!\left[
\theta_\xi:H^0(C,K_C\otimes W)\longrightarrow H^1(C,W)
\right].
\]
Then
\begin{equation}\label{eq:common-kernel-identification}
K_W^{1,0}\simeq H^0(C,Q^\vee\otimes K_C^2)
\end{equation}
and
\begin{equation}\label{eq:common-kernel-bound}
\dim K_W^{1,0}
\ge\max\{0,c(3g-1)-2r\}.
\end{equation}
\end{proposition}

\begin{proof}
Under Serre duality, the adjoint of the infinitesimal period map is trace multiplication
\[
H^0(K_C\otimes W)\otimes H^0(K_C\otimes W^\vee)
\longrightarrow H^0(K_C^2).
\]
More explicitly, for $s\in H^0(K_C\otimes W)$, $t\in H^0(K_C\otimes W^\vee)$, and $\xi\in H^1(T_C)$,
\[
\langle\theta_\xi(s),t\rangle_{\mathrm{SD}}
=\langle\xi,\operatorname{tr}(s\otimes t)\rangle_{\mathrm{SD}}.
\]
Thus a section in the first factor is killed for every tangent vector precisely when the corresponding homomorphism
$F=K_C\otimes W^\vee\to K_C^2$ vanishes on the evaluation image.  It then vanishes on its saturation $U$, because a torsion sheaf has no nonzero map to the line bundle $K_C^2$.  It therefore factors uniquely through $Q$, proving \eqref{eq:common-kernel-identification}; the converse follows from the same Serre-duality identity.

Since $\rk Q=c$ and $\deg Q=2r(g-1)-D$, Riemann--Roch gives
\[
\chi(Q^\vee\otimes K_C^2)
=D+(3c-2r)(g-1).
\]
Using \eqref{eq:Dlower}, this is at least $c(3g-1)-2r$.  Since $h^0\ge\chi$, equation \eqref{eq:common-kernel-bound} follows.
\end{proof}

\begin{remark}\label{rem:kernel-not-flat}
Proposition~\ref{prop:common-higgs} is a fibrewise first-order statement.  It neither makes $K_W^{1,0}$ Gauss--Manin horizontal nor gives it parabolic degree zero or a rational structure.  A Putman--Wieland finite orbit would instead produce, after finite level, a nonzero rational flat sublocal system; in Higgs terms, a degree-zero horizontal summand.  The proposition supplies none of these global properties.
\end{remark}

\section{Finite orbits and maximal-rank Prym maps}\label{sec:pw}

I first record the exact dictionary between finite monodromy orbits and constant Hodge-theoretic factors.  Let $S$ be a connected smooth complex quasi-projective variety, let $\mathbb V_{\Z}$ be a polarizable integral variation of Hodge structure on $S$, and write $\Gamma=\pi_1(S,s)$.

\begin{proposition}[Finite orbit dictionary]\label{prop:finite-orbit-dictionary}
The following conditions are equivalent:
\begin{enumerate}[label=\textup{(\roman*)}]
\item $\Gamma$ has a nonzero finite orbit in $V_{s,\C}$;
\item after a connected finite \'etale base change $T\to S$, the pulled-back local system has a nonzero global flat section;
\item after a connected finite \'etale base change, $\mathbb V_{\Q}$ contains a nonzero constant $\Q$-subvariation of Hodge structure.
\end{enumerate}
\end{proposition}

\begin{proof}
A vector has finite orbit exactly when its stabilizer has finite index.  By the Riemann existence theorem, the corresponding finite topological cover of $S$ is algebraic and finite \'etale; taking the connected component determined by the stabilizer gives (ii), and the converse is immediate.

If a finite-index subgroup $\Delta$ fixes a nonzero complex vector, then
\[
V_{\C}^{\Delta}=V_{\Q}^{\Delta}\otimes_{\Q}\C\ne0,
\]
because invariants of a rational representation commute with scalar extension.  The theorem of the fixed part makes the maximal trivial sublocal system of a polarizable variation a constant subvariation of Hodge structure \cite[Section~6(a)]{GriffithsSchmid1975}.  This proves (ii)$\Rightarrow$(iii); the reverse implication is immediate.
\end{proof}

The same dictionary holds for a smooth quasi-projective Deligne--Mumford stack admitting a finite \'etale surjective cover by a smooth quasi-projective variety: pull any alleged finite orbit or constant subvariation to a connected component of such a cover and apply Proposition~\ref{prop:finite-orbit-dictionary} there.

\begin{theorem}[Maximal-rank fixed-part obstruction]\label{thm:no-fixed-part}
Let $S$ be a connected smooth complex quasi-projective variety and let $\mathbb V_{\Z}$ be a polarizable integral variation of Hodge structure of type $\{(1,0),(0,1)\}$ and rank $2n$.  Suppose that at some $s\in S$ the differential of its local period map has rank
\begin{equation}\label{eq:full-siegel-dimension}
\rk(d\Phi_s)=\frac{n(n+1)}2.
\end{equation}
Then, for every connected finite \'etale map $q:T\to S$, the variation $q^*\mathbb V_{\Q}$ contains no nonzero constant subvariation.  Equivalently, neither the monodromy of $\mathbb V_{\C}$ nor that of any such pullback has a nonzero finite orbit.

The same statement holds for a smooth Deligne--Mumford stack admitting a finite \'etale surjective cover by a smooth quasi-projective variety.
\end{theorem}

\begin{proof}
Suppose that a constant subvariation
$0\ne\mathbb U_{\Q}\subset q^*\mathbb V_{\Q}$ exists.  Polarizability gives an orthogonal decomposition
\[
q^*\mathbb V_{\Q}=\mathbb U_{\Q}\oplus\mathbb W_{\Q}.
\]
Write $\rk_{\Q}\mathbb U=2a$, where $1\le a\le n$.  On Hodge bundles,
\[
E^{1,0}=U^{1,0}\oplus W^{1,0},
\qquad \rk U^{1,0}=a,
\qquad \rk W^{1,0}=n-a.
\]
The constant summand has zero Higgs field.  The two summands are flat and orthogonal, so there are no off-diagonal Higgs terms.  Consequently the differential of the period map takes values in
\[
\Sym^2\bigl((W^{1,0})^\vee\bigr)
\subset\Sym^2\bigl((E^{1,0})^\vee\bigr),
\]
and at every point of $T$,
\begin{equation}\label{eq:period-rank-defect}
\rk d(\Phi\circ q)
\le\frac{(n-a)(n-a+1)}2
<\frac{n(n+1)}2.
\end{equation}
Choose $t\in T$ above the point $s$ in \eqref{eq:full-siegel-dimension}.  Since $q$ is \'etale, $d(\Phi\circ q)_t$ has rank $n(n+1)/2$, contradicting \eqref{eq:period-rank-defect}.  Proposition~\ref{prop:finite-orbit-dictionary}, applied also after further finite covers, gives the finite-orbit formulation.  For a Deligne--Mumford stack, pass to a smooth finite level cover and apply the variety case.
\end{proof}

In particular, if $\dim S=n(n+1)/2$ and the classifying morphism
$S\to\cA_{n,\delta}$ is dominant and generically finite, then the hypothesis holds on a nonempty open subset.  Maximal rank, rather than generic finiteness onto an arbitrary image, is essential: the map $B\mapsto A_0\times B$ may be generically finite onto a product locus while the associated variation contains the constant factor $H^1(A_0,\Q)$.

\subsection{The degree-\texorpdfstring{$27$}{27} Prym map}

Let $\cR_6$ be the stack of connected \'etale double covers
\[
\widetilde C\longrightarrow C,
\qquad g(C)=6,
\]
and let
\[
P_6:\cR_6\longrightarrow\cA_5
\]
be the Prym map.  The anti-invariant local system
\[
\mathbb V_6^-=(R^1\widetilde f_*\Q)^-
\]
has rank ten and is the weight-one variation classified by $P_6$.  One has
\[
\dim\cR_6=15=\dim\cA_5,
\]
and Donagi--Smith proved that $P_6$ is dominant and generically finite of degree $27$ \cite[Theorem~I.2.1]{DonagiSmith1981}.  Adding full level $N\ge3$ to the curve and Prym lattice gives a smooth quasi-projective finite \'etale cover, so Theorem~\ref{thm:no-fixed-part} applies to this stack.

\begin{corollary}[Virtual Prym fixed parts in genus six]\label{cor:P6}
For every connected finite \'etale base change $\mathscr T\to\cR_6$, the variation $\mathbb V_6^-|_{\mathscr T}$ has no nonzero constant rational subvariation.  Equivalently, the anti-invariant Prym monodromy has no nonzero finite orbit.

For a fixed connected \'etale double cover of a genus-six surface, the full liftable mapping-class-group action on $H^1$ has no nonzero finite orbit.  Thus the cover satisfies the Putman--Wieland no-finite-orbit conclusion in the liftable-stabilizer formulation.
\end{corollary}

\begin{proof}
The first assertion is Theorem~\ref{thm:no-fixed-part}.  For the second, the rational deck decomposition is
\[
H^1(\widetilde C,\Q)=H^1(\widetilde C,\Q)^+
\oplus H^1(\widetilde C,\Q)^-.
\]
The minus part is covered by the first assertion.  The plus part is pulled back from $H^1(C,\Q)$.  The stabilizer of the defining nonzero two-torsion class is finite index in the mapping class group, and its image in $\Sp_{12}(\Z)$ is finite index.  The standard symplectic representation has no nonzero vector fixed by a finite-index subgroup.  Thus neither summand has a nonzero finite orbit.
\end{proof}

This is not a new case of the Putman--Wieland conjecture: it already follows from \cite[Theorem~1.2]{LandesmanLitt2023}, because the deck group has order $2<6^2$.  The result above is an independent maximal-period-rank proof and will be used for the no-extension interpretation of special Prym fibres.

At $(C,\eta)$, where $0\ne\eta\in\Pic^0(C)[2]$, the codifferential is
\begin{equation}\label{eq:P6-codifferential}
(dP_6)^*:\Sym^2H^0(C,K_C\otimes\eta)
\longrightarrow H^0(C,K_C^2).
\end{equation}
Both sides have dimension $15$, and the degree theorem implies that \eqref{eq:P6-codifferential} is an isomorphism on a dense open subset; see also \cite[Proposition~7.5]{Beauville1977}.

\begin{corollary}[Special Prym flat sections do not extend]\label{cor:no-thickening}
Let $Z\to\cR_6$ be a connected smooth family contained in a fibre of $P_6$, for example a smooth-Hurwitz open of the Fano-surface fibre over the intermediate Jacobian of a cubic threefold \cite[Theorem~4.1]{CasalainaMartinZhang2021}.  Let $\mathscr T\to\cR_6$ be connected and finite \'etale, and let $Z'$ be a connected component of $Z\times_{\cR_6}\mathscr T$.  No nonzero flat section of $\mathbb V_6^-|_{Z'}$ extends to a global flat section of $\mathbb V_6^-|_{\mathscr T}$.
\end{corollary}

\begin{proof}
An extension would be a nonzero monodromy-invariant vector on $\mathscr T$, contrary to Corollary~\ref{cor:P6}.
\end{proof}

The proper allowable-cover fibre is the Fano surface, whereas its intersection with the smooth Hurwitz stack is only a Zariski-open part.  Corollary~\ref{cor:no-thickening} makes no assertion about subvariations supported entirely on the boundary.

\subsection{The ramified degree-\texorpdfstring{$3$}{3} analogue}

Let $\cR_{3,4}$ parametrize triples $(C,\eta,D)$ with $g(C)=3$, $D$ reduced of degree four, and
$\eta^{\otimes2}\simeq\OO_C(D)$.  The corresponding double cover is ramified at four points.  Its Prym has dimension four and polarization type $\delta=(1,2,2,2)$, up to polarized duality.  The Prym map
\[
P_{3,4}:\cR_{3,4}\longrightarrow\cA_{4,\delta}
\]
has equal source and target dimensions $10$ and is dominant of generic degree three \cite{NagarajRamanan1995,BardelliCilibertoVerra1995}; see also \cite[Corollary~0.3]{CasalainaMartinZhang2021}.  Adding full level $N\ge3$ gives a smooth quasi-projective finite \'etale cover, so Theorem~\ref{thm:no-fixed-part} applies.

\begin{corollary}[Ramified Prym fixed parts]\label{cor:P34}
After no connected finite \'etale base change of $\cR_{3,4}$ does its anti-invariant Prym variation contain a nonzero constant rational subvariation; equivalently, its monodromy has no nonzero finite orbit.

For a fixed branched double cover $\widetilde C\to C$ in this Hurwitz component, the full liftable-stabilizer action on $H^1(\widetilde C,\Q)$ has no nonzero finite orbit.

Let $E^\circ$ be the smooth-Hurwitz open of the special elliptic fibre in \cite[Theorem~5.1]{CasalainaMartinZhang2021}.  If $\mathscr T\to\cR_{3,4}$ is connected and finite \'etale and $E'$ is a connected component of $E^\circ\times_{\cR_{3,4}}\mathscr T$, then no nonzero flat anti-invariant Prym section on $E'$ extends to a global flat section on $\mathscr T$.
\end{corollary}

\begin{proof}
The first assertion follows from maximal generic rank and Theorem~\ref{thm:no-fixed-part}.  For the full cohomology, use the rational deck decomposition
\[
H^1(\widetilde C,\Q)=H^1(\widetilde C,\Q)^+
\oplus H^1(\widetilde C,\Q)^-.
\]
The minus part is the Prym variation.  The plus part is $H^1(C,\Q)$.  The forgetful map from the pure four-punctured mapping class group onto the genus-three mapping class group is surjective; therefore the image of the finite-index liftable stabilizer is finite index, as is its image in $\Sp_6(\Z)$.  The standard representation has no nonzero finite orbit.  Finally, an extension of a section from $E'$ would contradict the first assertion.
\end{proof}

This double-cover no-finite-orbit case was already known from \cite[Theorem~1.2]{LandesmanLitt2023}, since the deck group has order $2<3^2$.  The corollary records the independent Prym-period proof and its consequence for the special smooth-Hurwitz fibre; it makes no assertion about boundary-supported subvariations in the allowable-cover compactification.

The codifferential at $(C,\eta,D)$ is
\begin{equation}\label{eq:P34-codifferential}
(dP_{3,4})^*:
\Sym^2H^0(C,K_C\otimes\eta)
\longrightarrow H^0(C,K_C^2(D)).
\end{equation}
Both spaces have dimension ten, so \eqref{eq:P34-codifferential} is an isomorphism at a general point; see \cite[Section~5.4]{CasalainaMartinZhang2021}.  Removing the four branch points gives an unramified double cover of a four-punctured genus-three surface, and filling the punctures upstairs recovers the original genus-seven covering curve.  Thus the anti-invariant variation is a punctured-surface higher-Prym variation in the sense of Putman--Wieland.  It is not the variation of an unramified cover of the closed genus-three surface.  The statement concerns the finite-index liftable stabilizer, not necessarily a characteristic subgroup acted on by the entire mapping class group.

\subsection{The exact boundary of the comparison}

For the geometric higher-Prym local system on a connected Hurwitz component, a Putman--Wieland finite orbit, a flat section after finite level, and a constant rational subvariation are equivalent by Proposition~\ref{prop:finite-orbit-dictionary}.  For the curve-cover variation considered by Litt, \cite[Proposition~6.3.2]{Litt2024} says that generic global generation precludes such a virtual section.  Its contraposition gives the one-way implication
\begin{equation}\label{eq:logical-chain}
\begin{aligned}
\text{finite orbit}
&\ \Longrightarrow\ 
\text{nonzero rational flat section after finite level}\\
&\ \Longrightarrow\
\text{failure of generic global generation}.
\end{aligned}
\end{equation}
Theorem~\ref{thm:uniform} establishes the last condition, but the implication cannot be reversed: it produces neither condition to its left.  Its deck group satisfies $|G_g|\ge1+g^{2g}>g^2$, so \cite[Theorem~1.2]{LandesmanLitt2023} does not settle this particular higher-Prym representation.  Conversely, Corollaries~\ref{cor:P6} and \ref{cor:P34} exclude the middle condition for the full Prym variations despite the presence of special positive-dimensional fibres.

The $26$-dimensional summand in Corollary~\ref{cor:26summand} is a weight-two variation on Fano surfaces, defined over a cyclotomic character field.  A Putman--Wieland finite orbit in complex weight-one cohomology forces, after finite level, the existence of a nonzero rational flat class.  The rank-seven period calculation proves that the distinguished complex component of the $F_4$ summand is not isotrivial along the marked Eckardt branch; it does not change its weight or supply the missing rational flat class.  This is the precise relationship established here between the $F_4$, Raynaud--Prill, and Putman--Wieland constructions.

\appendix

\section{Exact Jacobian-ring certificates}\label{app:certificates}

All entries below are obtained in the Artinian ring \eqref{eq:jacobian-ring}; hence the certificates are computations over $\Q$.  They are included in full so that the rank assertions do not depend on numerical linear algebra.  Products $a_i a_j$ and $u_k u_l$ denote the unnormalized products in the corresponding symmetric algebras; no factor $1/2$ is inserted.

\subsection{The rank-\texorpdfstring{$27$}{27} restricted cup product}

Use the basis
\[
R_4=\langle
x_0x_2x_3x_4,
x_1x_2x_3x_4,
x_1^2x_3x_4,
x_1^2x_2x_4,
x_1^2x_2x_3
\rangle
\]
and the socle generator $x_1^2x_2x_3x_4$.  The latter makes the perfect pairing
$R_4\otimes R_1\to R_5\simeq\Q$ explicit and hence identifies
$\wedge^2R_4$ with $(\wedge^2R_1)^\vee$.  This identification is used for the central and flattened matrices; the $27\times27$ certificate instead uses direct $\wedge^2R_4$-coordinates.

Order the standard basis of $\wedge^2R_1$, and the basis of $\wedge^2R_4$ induced by the displayed $R_4$-basis, by
\[
01,02,03,04,12,13,14,23,24,34.
\]
Thus a row $(ij,kl)$ records the $kl$-coordinate of
$\nu(\,cdot\,)(x_i\wedge x_j)$ in $\wedge^2R_4$.
Take all columns $a_i a_j$, $0\le i\le j\le6$, except $a_3a_4$, and the following rows of the matrix of $\nu$:
\begin{gather*}
(01,01),(01,02),(01,03),(01,04),(02,03),(02,04),(03,04),\\
(12,12),(12,13),(12,14),(12,23),(12,24),(12,34),\\
(13,13),(13,14),(13,23),(13,24),(13,34),\\
(14,14),(14,24),(14,34),\\
(23,23),(23,24),(23,34),(24,24),(24,34),(34,34).
\end{gather*}
The resulting $27\times27$ determinant is $-128$.  Direct substitution in \eqref{eq:nu} gives
\[
\nu(a_1a_6+a_2a_5+a_3a_4)=0.
\]
These two facts prove \eqref{eq:kernel-mu}.

\subsection{The central rank-\texorpdfstring{$12$}{12} minor}

With columns
\[
\begin{split}
&(u_0^2,u_0u_1,u_0u_2,u_1^2,u_1u_2,u_2^2;\\
&\hspace{32mm}u_3^2,u_3u_4,u_3u_5,u_4^2,u_4u_5,u_5^2)
\end{split}
\]
and rows
\[
(a_0^2,a_0a_1,a_0a_2,a_0a_4,a_0a_5,a_0a_6,
a_1^2,a_1a_2,a_1a_4,a_1a_6,a_2^2,a_4^2),
\]
where rows are $T_E$-products and columns are $U_\alpha$-products, the relevant minor is
\begin{equation}\label{eq:matrix12}
\begin{pmatrix}
-6&0&0&0&0&0&0&0&0&0&0&0\\
0&-3&0&0&0&0&0&0&0&0&0&0\\
0&0&-3&0&0&0&0&0&0&0&0&0\\
0&0&0&0&-3&0&0&1&0&0&0&0\\
0&3&0&6&3&0&-2&-1&0&0&0&0\\
0&0&3&0&3&6&0&-1&0&-2&0&0\\
0&0&0&0&0&0&-2&0&0&0&0&0\\
0&0&0&0&0&0&0&-1&0&0&0&0\\
0&0&0&0&0&0&0&0&1&0&0&0\\
0&0&0&0&0&0&0&1&-1&0&-2&0\\
0&0&0&0&0&0&0&0&0&-2&0&0\\
0&0&0&0&0&0&0&0&0&0&0&-2
\end{pmatrix},
\qquad
\det=-2^7 3^6.
\end{equation}

\subsection{The flattened rank-\texorpdfstring{$7$}{7} minor}

Take columns $a_0,\ldots,a_6$ and rows
\[
(a_0,u_0^2),(a_0,u_0u_1),(a_0,u_0u_2),(a_0,u_1^2),
(a_0,u_1u_2),(a_0,u_2^2),(a_1,u_0u_1).
\]
A row $(\eta,u_ku_l)$ and a column $\xi$ evaluate
$\overline\Theta^{(2)}_0(\xi\eta)$ on $u_ku_l$.  The minor is
\begin{equation}\label{eq:matrix7}
\begin{pmatrix}
-6&0&0&6&0&0&0\\
0&-3&0&3&0&3&0\\
0&0&-3&3&0&0&3\\
0&0&0&0&0&6&0\\
0&0&0&0&-3&3&3\\
0&0&0&0&0&0&6\\
-3&0&0&0&0&0&0
\end{pmatrix},
\qquad
\det=-2^3 3^7.
\end{equation}
An exact rational-arithmetic verifier, \texttt{verify\_exact\_minors.py}, accompanies the source.  The displayed nonzero determinants themselves are the proof certificates; the script only reproduces them.

\section*{Acknowledgments}

GPT-5.6 sol was used for exploratory computation, proof checking, and language editing.  Responsibility for the statements and arguments remains with the author.

\begin{thebibliography}{99}

\bibitem{BPGN1997}
L.~Brambila-Paz, I.~Grzegorczyk, and P.~E. Newstead,
\emph{Geography of Brill--Noether loci for small slopes},
J. Algebraic Geom. \textbf{6} (1997), 645--669.

\bibitem{BardelliCilibertoVerra1995}
F.~Bardelli, C.~Ciliberto, and A.~Verra,
\emph{Curves of minimal genus on a general abelian variety},
Compositio Math. \textbf{96} (1995), 115--147.

\bibitem{Beauville1977}
A.~Beauville,
\emph{Vari\'et\'es de Prym et jacobiennes interm\'ediaires},
Ann. Sci. \`Ecole Norm. Sup. (4) \textbf{10} (1977), 309--391.

\bibitem{BBD1982}
A.~Beilinson, J.~Bernstein, and P.~Deligne,
\emph{Faisceaux pervers},
Ast\'erisque \textbf{100} (1982), 5--171.

\bibitem{CasalainaMartinZhang2021}
S.~Casalaina-Martin and Z.~Zhang,
\emph{The moduli space of cubic surface pairs via the intermediate Jacobians of Eckardt cubic threefolds},
J. Lond. Math. Soc. (2) \textbf{104} (2021), no.~1, 1--34.

\bibitem{ClemensGriffiths1972}
C.~H. Clemens and P.~A. Griffiths,
\emph{The intermediate Jacobian of the cubic threefold},
Ann. of Math. (2) \textbf{95} (1972), 281--356.

\bibitem{deCataldoMigliorini2002}
M.~A.~A. de~Cataldo and L.~Migliorini,
\emph{The hard Lefschetz theorem and the topology of semismall maps},
Ann. Sci. \`Ecole Norm. Sup. (4) \textbf{35} (2002), 759--772.

\bibitem{DonagiSmith1981}
R.~Donagi and R.~C. Smith,
\emph{The structure of the Prym map},
Acta Math. \textbf{146} (1981), 25--102.

\bibitem{GriffithsSchmid1975}
P.~Griffiths and W.~Schmid,
\emph{Recent developments in Hodge theory: a discussion of techniques and results},
in \emph{Discrete Subgroups of Lie Groups and Applications to Moduli},
Oxford Univ. Press, 1975, 31--127.

\bibitem{JKLM2025}
A.~Javanpeykar, T.~Kr\"amer, C.~Lehn, and M.~Maculan,
\emph{The monodromy of families of subvarieties on abelian varieties},
Duke Math. J. \textbf{174} (2025), 1045--1149.

\bibitem{KLMExceptional2026}
T.~Kr\"amer, C.~Lehn, and M.~Maculan,
\emph{Exceptional Tannaka groups only arise from cubic threefolds},
J. Reine Angew. Math. \textbf{830} (2026), 51--99.

\bibitem{KLM2026}
T.~Kr\"amer, D.~Litt, and M.~Maculan,
\emph{$E_6$-local systems from cubic threefolds},
arXiv:2604.20970 (2026).

\bibitem{Kramer2016}
T.~Kr\"amer,
\emph{Cubic threefolds, Fano surfaces and the monodromy of the Gauss map},
Manuscripta Math. \textbf{149} (2016), 303--314.

\bibitem{Kramer2022}
T.~Kr\"amer,
\emph{Characteristic cycles and the microlocal geometry of the Gauss map, I},
Ann. Sci. \`Ecole Norm. Sup. (4) \textbf{55} (2022), 1475--1527.

\bibitem{KramerWeissauer2015}
T.~Kr\"amer and R.~Weissauer,
\emph{Vanishing theorems for constructible sheaves on abelian varieties},
J. Algebraic Geom. \textbf{24} (2015), 531--568.

\bibitem{LandesmanLitt2023}
A.~Landesman and D.~Litt,
\emph{An introduction to the algebraic geometry of the Putman--Wieland conjecture},
European J. Math. \textbf{9} (2023), Article~40.

\bibitem{LiebeckSeitz2004}
M.~W. Liebeck and G.~M. Seitz,
\emph{Subgroups of exceptional algebraic groups which are irreducible on an adjoint or minimal module},
J. Group Theory \textbf{7} (2004), 347--372.

\bibitem{Litt2024}
D.~Litt,
\emph{Motives, mapping class groups, and monodromy},
arXiv:2409.02234v1 (2024).

\bibitem{Looijenga1997}
E.~Looijenga,
\emph{Prym representations of mapping class groups},
Geom. Dedicata \textbf{64} (1997), 69--83.

\bibitem{NagarajRamanan1995}
D.~S. Nagaraj and S.~Ramanan,
\emph{Polarisations of type $(1,2,\ldots,2)$ on abelian varieties},
Duke Math. J. \textbf{80} (1995), 157--194.

\bibitem{NarasimhanSeshadri1965}
M.~S. Narasimhan and C.~S. Seshadri,
\emph{Stable and unitary vector bundles on a compact Riemann surface},
Ann. of Math. (2) \textbf{82} (1965), 540--567.

\bibitem{PutmanWieland2013}
A.~Putman and B.~Wieland,
\emph{Abelian quotients of subgroups of the mapping class group and higher Prym representations},
J. Lond. Math. Soc. (2) \textbf{88} (2013), 79--96.

\bibitem{Raynaud1982}
M.~Raynaud,
\emph{Sections des fibr\'es vectoriels sur une courbe},
Bull. Soc. Math. France \textbf{110} (1982), 103--125.

\bibitem{Roulleau2009}
X.~Roulleau,
\emph{Elliptic curve configurations on Fano surfaces},
Manuscripta Math. \textbf{129} (2009), 381--399.

\bibitem{Schneider2007}
O.~Schneider,
\emph{Sur la dimension de l'ensemble des points base du fibr\'e d\'eterminant sur $\mathcal{SU}_C(r)$},
Ann. Inst. Fourier (Grenoble) \textbf{57} (2007), 481--490.

\bibitem{SpringerVeldkamp2000}
T.~A. Springer and F.~D. Veldkamp,
\emph{Octonions, Jordan Algebras and Exceptional Groups},
Springer Monographs in Mathematics, Springer, Berlin, 2000.

\end{thebibliography}

\end{document}
